\clearpage
\VerdierPhysicalPageAnchor{0203}%
\begingroup\small\setlength{\parskip}{0pt}\setlength{\emergencystretch}{1.5em}
\setlength{\abovedisplayskip}{4pt plus 2pt minus 2pt}%
\setlength{\belowdisplayskip}{4pt plus 2pt minus 2pt}%

\VerdierSemanticAnchor{ch3d-2-4-1-cont}
\noindent\textit{は圏同値である。自然関手}
\[
K^{-}(M^{\wedge})/\operatorname{Ac}^{-}(M^{\wedge})
\longrightarrow D^{-}(\mathcal A)
\]
\noindent\textit{は充満忠実である。コホモロジー対象が $M^{\wedge}$ に属する
$D^{-}(\mathcal A)$ の対象は、
$K^{-}(M^{\wedge})/\operatorname{Ac}^{-}(M^{\wedge})$ の対象と同型である。
ここで $\operatorname{Ac}^{-}(M)$ と $\operatorname{Ac}^{-}(M^{\wedge})$ は、
その成分がそれぞれ $M$ と $M^{\wedge}$ に属する非輪状複体の圏を表す。}

\VerdierSemanticAnchor{ch3d-2-4-1-proof}
この命題は (2.3.1)、(2.3.2)、（第 II 章，2.3.5）から直ちに従う。

\VerdierSemanticAnchor{ch3d-2-4-2}
\medskip\noindent\textbf{系 2.4.2.}\quad
\textit{$U$ を宇宙とする。$\mathcal A$ は $U$-圏、すなわち
$\mathcal A$ の任意の二対象 $Y$、$Z$ に対して
$\operatorname{Hom}_{\mathcal A}(Y,Z)$ は $U$ の元であり、さらに $U$ の元で
添字づけられた直和が $\mathcal A$ において表現可能であると仮定する。
$(X_i)_{i\in I\in U}$ を $\mathcal A$ の生成元の族とし、$M$ を
$\mathcal A$ の対象のうち $\coprod_{j\in J\in U}X_j$ の形をするもの全体とする。
ここで各 $X_j$ は集合 $\bigcup_{i\in I}\{X_i\}$ の元である。
このとき自然関手}
\[
K^{-}(M)/\operatorname{Ac}^{-}(M)\longrightarrow D^{-}(\mathcal A)
\]
\noindent\textit{は圏同値である。}

\VerdierSemanticAnchor{ch3d-2-4-2-proof}
これは (2.4.1) から直ちに従う。

\VerdierSemanticAnchor{ch3d-2-4-3}
\medskip\noindent\textbf{命題 2.4.3.}\quad
\textit{$\mathcal A$ をアーベル圏とし、$M$ を性質 $(E_4)$ と $(E_5)$ をもつ
$\mathcal A$ の対象の集合とする。$\mathcal A$ の各対象は $M$ のある元の商対象であり、
さらに $\mathcal A$ の各対象 $X$ の $M$-次元を上から抑える整数
$n\in\mathbb N$ が存在すると仮定する。}

\noindent\textit{a) 圏 $K^{*}(M)$ は左
$\operatorname{Ac}^{*}(\mathcal A)$-局所化部分圏である
（$*=-,+,b,{}$「空」）。}

\noindent\textit{b) 自然関手}
\[
K^{*}(M)/\operatorname{Ac}^{*}(M)\longrightarrow D^{*}(\mathcal A)
\]
\noindent\textit{は圏同値である。}

\VerdierSemanticAnchor{ch3d-2-4-3-proof}
これは命題 (2.3.3)（$*=+,b,{}$「空」）、または命題 (2.3.1)（$*=-$）の
言い換えである。

\VerdierSemanticAnchor{ch3d-2-4-4}
\medskip\noindent\textbf{命題 2.4.4.}\quad
\textit{$Q^{*}\colon K^{*}(\mathcal A)\longrightarrow D^{*}(\mathcal A)$ を
標準関手とする。}

\noindent\textit{a) $X$ を $\mathcal A$ の対象とする。$X$ が $Q$-左自由
（第 II 章，2.3.2）であるための必要十分条件は、$X$ が圏 $\mathcal A$ の
射影対象であることである。}
\endgroup
